\section{Conclusions and Future Work}
In this work, we introduced PromptCoT~2.0, a principled framework for rationale-driven prompt synthesis that addresses the persistent shortage of high-quality training problems for reasoning-focused LLMs.  
By formulating concept–rationale–prompt generation as an expectation–maximization procedure, our method iteratively refines rationales and leverages them to construct substantially more difficult and diverse problems across mathematics and programming.  
Beyond synthesis, we demonstrated how these problems can fuel two complementary post-training regimes: \emph{self-play}, which enables strong models to improve autonomously through verifiable feedback without reliance on ever-stronger external teachers, and \emph{supervised fine-tuning}, which allows weaker models to acquire reasoning skills from teacher-distilled traces.  Extensive experiments show that PromptCoT~2.0 not only establishes new state-of-the-art results at the 30B scale but also allows 7B models trained purely on synthetic prompts to achieve performance competitive with models trained on human-curated or hybrid datasets. 
Our analyses further confirm that the synthesized problems differ fundamentally from prior corpora, exhibiting higher difficulty and richer distributional diversity—qualities essential for advancing LLM reasoning.  

Looking ahead, we view prompt synthesis as a central axis for scaling reasoning beyond brute-force model size or computation.  
Future work includes extending our EM-based framework to multimodal settings, integrating richer verification signals to broaden self-play, and exploring its role in the development of agentic intelligence.  
We hope that PromptCoT~2.0 and the accompanying dataset will provide a foundation for the next generation of open-source reasoning models and accelerate progress toward autonomous, verifiably strong LLMs.
